
\section{闭嵌入与闭子概形}

\textbf{9.1.1.}\ \emph{\trhDefinition{}。}\quad 概形之间的\trMorphism{} $\pi\colon X\to Y$ 称为\textbf{\trClosedEmbedding}（closed embedding；或\textbf{\trClosedImmersion}，closed immersion），若它是一个\trAffine\trMorphism，并且对 $Y$ 的每个\trAffine\trOpenSubset{} $\Spec B\subset Y$，只要 $\pi^{-1}(\Spec B)\cong\Spec A$，映射 $B\to A$ 就都是\trSurjectiveMap{}（也就是说，形如 $B\to B/I$——我们心仪的那个局部模型）。（「嵌入」比「浸入」更可取，因为微分几何中的浸入观念，更接近代数几何学家所称的非分歧（unramified），我们将在 \S21.7 中给出它的定义。）符号 $\hookrightarrow$ 常被用来表示一个态射是\trClosedEmbedding{}（或更一般地，是\trLocally\trClosedEmbedding，见 \S9.2）。

若 $X$ 是 $Y$ 的一个\emph{子集}（且 $\pi$ 在集合的层面上就是包含映射），我们就说 $X$ 是 $Y$ 的一个\textbf{\trClosedSubscheme}。一个\trClosedEmbedding{}与「向一个\trClosedSubscheme{}的\trIsomorphism」是一回事。

\textbf{9.1.A.（练习）}\quad 证明：一个\trClosedEmbedding{}把 $X$ 的底拓扑空间与 $Y$ 的底拓扑空间的一个\trClosedSubset{}等同起来。（\trhCaveat{}：两道\trClosedEmbedding{} $\Spec k[x]/(x)\hookrightarrow\Spec k[x]$ 与 $\Spec k[x]/(x^2)\hookrightarrow\Spec k[x]$ 表明，\trClosedSubset{}并不能决定\trClosedSubscheme：那些「无穷小」的信息——或者说「绒毛」——丢失了。）

\textbf{9.1.B.（容易的练习）}\quad 证明：\trClosedEmbedding{}都是\trFinite\trMorphism，从而是\trFiniteType{}的。

\textbf{9.1.C.（容易的练习）}\quad 证明：两个\trClosedEmbedding{}的\trComposition{}仍是一个\trClosedEmbedding。

\textbf{9.1.D.（练习）}\quad 证明：「是\trClosedEmbedding」这一性质是对目标仿射局部的。

特别地，若 $B\to A$ 是环之间的一个\trSurjectiveMap，则由之诱导的\trMorphism{} $\Spec A\to\Spec B$ 是一个\trClosedEmbedding。

\textbf{9.1.2.\ 重要：\trClosedSubscheme{}对应于\trQC{}的\trIdealSheaf{}。}\quad 我们把 $\OO_Y$ 的一个 $\OO_Y$-\trSubmodule{}称为 $Y$ 上的一个\textbf{\trIdealSheaf{}}（ideal sheaf，或 sheaf of ideals）。%
一个\trClosedSubscheme{} $\pi\colon X\hookrightarrow Y$ 给出 $\OO_Y$-\trModule{}的一个\trSurjectiveMap{} $\OO_Y\to\pi_*\OO_X$。该映射的\trKernel{} $\scr{I}_{X/Y}$ 便是 $Y$ 上的一个\trIdealSheaf{}。在每个\trOpenSubset{} $\mathsf{U}$ 上，它给出环 $\OO_Y(\mathsf{U})$ 的一个理想 $\scr{I}_{X/Y}(\mathsf{U})$；而在仿射开子集 $\Spec B$ 上，它给出 $B$ 的理想 $I$，即「\trClosedSubscheme」的定义中决定满射 $B\to A=B/I$ 的那个理想。

\textbf{9.1.E.（练习）}\quad 证明：$\pi_*\OO_X$ 是 $Y$ 上的一个\trQCSheaf{}。（一条可行的进路是用\mbox{练习~6.2.D}。）

\trQCSheaf{}构成一个\trAbelianCategory，故 $\scr{I}_{X/Y}=\ker(\OO_Y\to\pi_*\OO_X)$ 也是 $Y$ 上的一个\trQCSheaf。于是，$Y$ 的每个\trClosedSubscheme{}都给出 $Y$ 上一个\emph{拟凝聚的}\trIdealSheaf{}。我们把（$Y$ 上诸拟凝聚层之间的）正合列
\begin{equation*}\tag{9.1.2.1}\label{eq:closed-subscheme-exact-sequence}
0\longrightarrow\scr{I}_{X/Y}\longrightarrow\OO_Y\longrightarrow\pi_*\OO_X\longrightarrow 0
\end{equation*}
称为 $\pi\colon X\hookrightarrow Y$ 的\textbf{闭子概形正合列}（closed subscheme exact sequence）。

我们可以经由 \eqref{eq:closed-subscheme-exact-sequence} 从\trIdealSheaf{} $\scr{I}_{X/Y}$ 把 $X$（连同它到 $Y$ 的映射 $\pi$）恢复出来：对 $Y$ 的每个仿射开子集 $\Spec B\subset Y$，在 $\Spec B$ 上取\trSection{}便给出正合列
\[
0\longrightarrow I\longrightarrow B\longrightarrow B/I\longrightarrow 0,
\]
而 $\Spec B/I$ 正是 $\pi^{-1}(\Spec B)$。

逆命题成立：

\textbf{9.1.F.（重要练习）}\quad 证明：$Y$ 上每一个拟凝聚的\trIdealSheaf{}都以这种方式产生 $Y$ 的一个\trClosedSubscheme。

这是一道很难的练习，所以作为提示，这里给出两种不同的做法；把二者以某种方式结合起来，或许正合你的心意。\emph{做法一}。对每个仿射开集 $\Spec B$，我们有\trClosedSubscheme{} $\Spec B/I(B)\hookrightarrow\Spec B$。(i) 对任意两个仿射开子概形 $\Spec A$ 与 $\Spec B$，证明：两个\trClosedSubscheme{} $\Spec A/I(A)\hookrightarrow\Spec A$ 与 $\Spec B/I(B)\hookrightarrow\Spec B$ 限制到二者的交上，给出\emph{同一个}\trClosedSubscheme。（提示：用在两个仿射开集中同时特异的开集去覆盖二者的交，命题~5.3.1。）于是，举例来说，我们便可以把这两个\trClosedSubscheme{}粘合起来，得到 $\Spec A\cup\Spec B$ 的一个\trClosedSubscheme。(ii) 用关于粘合概形的\cref{ex:glue-arbitrary-schemes}（或其中的想法），把一切仿射开子概形中的\trClosedEmbedding{}同时粘合起来。你只需要操心三重交。\emph{做法二}。(i) 先把 $X$ 描述为 $Y$ 的一个子集。(ii) 检验 $X$ 是闭的。(iii) 在这个子集上定义函数层 $\OO_X$，或许会用到\trCompatibleGerm。(iv) 检验：这样得到的\trRingedSpace{}在局部上确实就是 $\Spec B/I\hookrightarrow\Spec B$ 给出的那个\trClosedSubscheme。

\textbf{9.1.G.（不重要练习：一个作为 $\OO$-模、却不\trQC{}的\trIdealSheaf{}）}\quad 设 $X=\Spec k[x]_{(x)}$，即「仿射直线在原点处的芽」，它有两个点：一个\trClosedPoint{}与一个\trGenericPoint{} $\eta$。定义 $\scr{I}(X)=\{0\}\subset\OO_X(X)=k[x]_{(x)}$ 与 $\scr{I}(\eta)=k(x)=\OO_X(\eta)$。证明：这个\trIdealSheaf{}并不对应于任何\trClosedSubscheme。

\textbf{9.1.H.}\ \extrahard\ \textbf{（困难的（而且坦白说不重要的）练习；后文不再用到）}\quad 文献中，\trClosedEmbedding{}的惯常定义是这样一个\trMorphism{} $\pi\colon X\to Y$：$\pi$ 诱导 $X$ 的底拓扑空间到 $Y$ 的底拓扑空间之某个\trClosedSubset{}上的一个\trHomeomorphism，而且 $Y$ 上诸\trSheaf{}之间由之诱导的映射 $\pi^{\sharp}\colon\OO_Y\to\pi_*\OO_X$ 是满的。（这里说「满」，指的是诸\trStalk{}上的\trRingMorphism{}是\trSurjectiveMap。）证明：这个定义与上面给出的定义相合。（要证明我们这条涉及「在仿射开集的层面上满」的定义蕴含这个定义，你可以使用「层的态射的满性可以在适当选取的\trBase{}上检验」这一事实，\cref{ex:surjectivity-on-base}。）

至此，我们已在概形的国度里定义了开子集与闭子集的类似物。\footnote{译注：「概形的国度」（the land of schemes）是作者的惯用戏语，即「概形的世界」。}它们的定义比起经典拓扑背景中的定义来，稍微不那么「对称」：一个\trClosedSubscheme{}的「\trComplement」是唯一的一个\trOpenSubscheme，但一般地，对一个给定的\trOpenSubscheme，却有许多「互补的」\trClosedSubscheme。（我们很快就会定义其中「最好的」那一个：带\trReduced{}结构的那个，\S9.4.9。）

\textbf{9.1.3.\ 例子与性质。}

\vkexsync{I}
\begin{exercise}[重要练习]\label{ex:scheme-theoretic-union-intersection}
\begin{enumerate}[label=(\alph*)]
\item 类比闭子集，定义 $X$ 的\textbf{闭子概形的有限（\trSchemeTheoretic{}）并}（finite (scheme-theoretic) union of closed subschemes）这一观念，以及任意的（不必有限的）\textbf{闭子概形的（\trSchemeTheoretic{}）交}（(scheme-theoretic) intersection of closed subschemes）。（\mbox{练习~9.1.F} 或有帮助。）提示：若 $X$ 仿射，那么你大概会指望：对应于 $I_1$ 与 $I_2$ 的两个\trClosedSubscheme{}之并，应当是对应于 $I_1\cap I_2$ 或 $I_1I_2$ 的那个\trClosedSubscheme——但究竟是哪一个？我们会希望一个\trClosedSubscheme{}与自身之并\emph{就是}它自身，所以正确的选择是 $I_1\cap I_2$。
\item 描述 $\mathbb{A}^2$ 中 $V(y-x^2)$ 与 $V(y)$ 的\trSchemeTheoretic{}交。图~\ref{fig:scheme-theoretic-intersection} 是一幅图。（比如说，请非正式地解释：这如何对应于两条曲线在单个点处以重数 $2$ 相交——注意在你的答案里那个 $2$ 是如何显现的。或者换个问法：非既约性就这个交向你泄露了什么——既是它的「大小」，也是它的「方向」？）再描述它们的\trSchemeTheoretic{}并。
\item 证明：\trClosedSubscheme{}的有限并的底集合，就是诸底集合的有限并；对任意交亦然。
\item 描述 $\mathbb{A}^2$ 中 $V(y^2-x^2)$ 与 $V(y)$ 的\trSchemeTheoretic{}交。画一幅图。（你先前料想这个交的重数是 $1$ 还是 $2$？）由此证明：若 $X$、$Y$、$Z$ 是 $W$ 的\trClosedSubscheme，则一般地 $(X\cap Z)\cup(Y\cap Z)\neq(X\cup Y)\cap Z$。特别地，并非交与并的一切性质都能从集合搬到概形上来。
\end{enumerate}
\end{exercise}

\textbf{9.1.J.（练习）}\quad 证明：\trClosedEmbedding{}（像\trOpenEmbedding{}一样；见容易的\cref{ex:open-embeddings-are-monomorphisms}）是\trMonomorphism。

在许多情形，先验地只是闭集的那些\trLocus，常常会有作为\trClosedSubscheme{}的「\trEnrichment{}」版本（enriched versions）。%
其余的练习就是这方面的例子。（如 \S6.6.33 所描述，我们常常通过添加形容词「\trSchemeTheoretic{}的」（scheme-theoretic）来表明：我们考虑的是\trSchemeTheoretic{}的\trEnrichment{}（scheme-theoretic enrichments）——比如说，我们会说「\trSchemeTheoretic{}交」「\trSchemeTheoretic{}并」，等等。）

\vkexsync{K}
\begin{exercise}[重要练习／定义：消失概形（参看\cref{ex:vanishing-scheme-definition}）]\label{ex:vanishing-scheme-global-section}
\begin{enumerate}[label=(\alph*)]
\item 设 $Y$ 是一个\trScheme，而 $s\in\Gamma(Y,\OO_Y)$。定义由 $s$ \textbf{切出}的那个\trClosedSubscheme。我们称之为 $s$ 的\textbf{\trVanishingScheme} $V(s)$，因为它是我们早先的（集合论版本的）$V(s)$ 的\trSchemeTheoretic{}版本（\S3.4.1）。（提示：在仿射开集 $\Spec B$ 上，我们径取 $\Spec B/(s_B)$，其中 $s_B$ 是 $s$ 限制到 $\Spec B$ 上的\trSection。用\mbox{练习~9.1.F} 证明：这给出一个\trWellDefined{}的\trClosedSubscheme。）
\item 若 $u$ 是一个可逆函数，证明：$V(s)=V(su)$。
\item 若 $S$ 是一个由函数组成的集合，定义 $V(S)$。
\end{enumerate}
\end{exercise}

\textbf{9.1.4.}\ \emph{\trhDefinition{}：\trSchemeTheoretic{}支撑。}\quad 若 $M$ 是一个 $A$-\trModule，而 $\mathrm{m}\in M$，我们称 $\Spec A$ 的闭子概形 $\Spec(A/\operatorname{Ann}\mathrm{m})$ 为 $\mathrm{m}$ 的\textbf{\trSchemeTheoretic{}支撑}（scheme-theoretic support）。%
（这常被描述为「$\mathrm{m}$ 在其上有定义的最小\trClosedSubscheme」。你能把这句话讲精确吗？）

\textbf{9.1.L.（练习）}\quad 证明：$\mathrm{m}$ 的\trSchemeTheoretic{}支撑的底集合，就是通常的「集合论」支撑 $\Supp\mathrm{m}$。

\textbf{9.1.M.（练习）}\quad 把有限生成 $A$-模 $M$ 的\textbf{\trSchemeTheoretic{}支撑}（scheme-theoretic support of a finitely generated $A$-module）定义为：使 $M$ 成为一个 $(A/I)$-\trModule{}的\emph{一切}\trClosedSubscheme{} $\Spec A/I\hookrightarrow\Spec A$ 的\trSchemeTheoretic{}交。证明：它等于 $M$ 的任意一组有限生成元中诸元的\trSchemeTheoretic{}支撑的\trSchemeTheoretic{}并。

\textbf{9.1.N.（练习：把上一道练习整体化）}\quad 设 $X$ 是一个\trScheme，而 $\scr{F}$ 是 $X$ 上一个\trFiniteType{}的\trQCSheaf。你将如何定义 $\scr{F}$ 的\textbf{\trSchemeTheoretic{}支撑}？
